% 闭凸集 C=[0,1]^2 上的投影与法锥。
% z_a=(1.6,0.6) 投影到边上的 p_a=(1,0.6)，法锥是水平射线；z_b=(1.7,-0.5) 投影到顶点 p_b=(1,0)，
% 法锥是象限 {v: v_1>=0, v_2<=0}。y=(0.3,0.9) 为 C 中任取的一点，<z_a-p_a, y-p_a> = 0.6*(-0.7) < 0。
\begin{tikzpicture}[scale=2.6,note/.style={font=\FigureNoteFont\fontsize{7.5}{9}\selectfont}]
  % normal cone at the vertex (1,0)
  \fill[FigureYellowFill] (1,0)--(1.95,0)--(1.95,-0.75)--(1,-0.75)--cycle;
  \draw[FigureYellow,line width=.6pt] (1.95,0)--(1,0)--(1,-0.75);
  \node[note,text=FigureYellow!70!FigureInk,anchor=west] at (1.42,-0.62) {顶点处的法锥};
  % the convex set
  \draw[FigureBlue,fill=FigureBlueFill,line width=.9pt] (0,0) rectangle (1,1);
  \node[text=FigureBlue] at (0.18,0.15) {$C$};
  % normal ray at the edge point
  \draw[FigureGreen,line width=.6pt,dashed] (1,0.6)--(1.95,0.6);
  % projection of z_a
  \draw[->,FigureGreen,line width=.9pt] (1,0.6)--(1.6,0.6);
  \fill[FigureInk] (1.6,0.6) circle (0.022) node[above right=-1pt] {$\bm z_a$};
  \fill[FigureGreen] (1,0.6) circle (0.022) node[above left=-1pt,text=FigureInk] {$\bm p_a$};
  % chord to an arbitrary feasible point
  \draw[->,FigureBlue,line width=.8pt] (1,0.6)--(0.3,0.9);
  \fill[FigureBlue] (0.3,0.9) circle (0.022) node[above=1pt,text=FigureInk] {$\bm y$};
  \draw[FigureMuted,line width=.5pt] (1.1,0.6) arc[start angle=0,end angle=156.8,radius=0.1];
  \node[note,text=FigureMuted] at (1.17,0.79) {钝角};
  % projection of z_b to the vertex
  \draw[->,FigureRed,line width=.9pt] (1,0)--(1.7,-0.5);
  \fill[FigureInk] (1.7,-0.5) circle (0.022) node[right=1pt] {$\bm z_b$};
  \fill[FigureRed] (1,0) circle (0.022) node[below left=-1pt,text=FigureInk] {$\bm p_b$};
  \node[note,text=FigureGreen!80!FigureInk,anchor=west] at (1.62,0.48) {边上的法锥};
\end{tikzpicture}
